% ================================================================ ANEXOS
\renewcommand{\thesection}{Anexo \Alph{section}}
\renewcommand{\thesubsection}{\Alph{section}.\arabic{subsection}}
\setcounter{section}{0}
\setcounter{table}{0}\renewcommand{\thetable}{A\arabic{table}}
\setcounter{figure}{0}\renewcommand{\thefigure}{A\arabic{figure}}

\begin{center}
{\Large\bfseries Anexos}\\[2pt]
{\small (material complementario; máximo 20 páginas)}
\end{center}

% ------------------------------------------------------------------ A
\section{Repositorio GitHub con los códigos fuente}
\label{anexo:repo}

Todo el análisis es reproducible desde el repositorio del grupo:

\begin{center}
\fbox{\parbox{.86\linewidth}{\centering\ttfamily
https://github.com/$\langle$usuario$\rangle$/proyecto1-btm-demanda\\[2pt]
\footnotesize (actualizar con el enlace definitivo al publicar; ver instrucciones abajo)}}
\end{center}

\subsection{Estructura}
\begin{verbatim}
proyecto1-btm-demanda/
|- README.md                  <- documentacion y guia de reproduccion
|- requirements.txt           <- numpy, pandas, statsmodels, matplotlib
|- data/
|  |- demanda_pib_anual.csv   <- demanda SIN (GWh) y PIB real, 2000-2025
|  |- upme_proyeccion_jul2024.csv <- proyeccion oficial UPME (esc. medio)
|  |- espana_anual.csv        <- demanda b.c., autoconsumo y PIB de Espana
|  |- parametros_btm.csv      <- parametros con fuente (CREG, APPA, PNUMA...)
|- src/
|  |- config.py               <- rutas, paleta de figuras y supuestos
|  |- analisis_econometrico.py<- ADF/KPSS, Johansen, Engle-Granger,
|  |                             VECM, VAR, combinacion de pronosticos
|  |- ajuste_btm.py           <- demanda bruta y escenarios BtM E0-E3
|  |- validacion.py           <- V1 backtest, V2 experimento Espana,
|  |                             V3 panel LSDV Colombia-Espana
|  |- generar_figuras.py      <- figuras F1-F7 del informe
|  |- descargar_datos.py      <- refresco opcional desde XM (pydataxm),
|  |                             Banco Mundial y anexo Excel UPME 2026-2040
|- output/                    <- tablas CSV y figuras (se regeneran)
|- informe/                   <- fuentes LaTeX de este documento
\end{verbatim}

\subsection{Reproducción}
\begin{verbatim}
pip install -r requirements.txt
cd src
python analisis_econometrico.py   # econometria y proyeccion organica
python ajuste_btm.py              # escenarios BtM y demanda neta
python validacion.py              # V1 backtest, V2 Espana, V3 panel
python generar_figuras.py         # figuras
\end{verbatim}
Los cuatro scripts corren en menos de un minuto y regeneran todas las tablas
y figuras de este informe. Para publicar el repositorio desde la cuenta de un
integrante: \texttt{git init \&\& git add .\ \&\& git commit}, crear el
repositorio vacío en GitHub y \texttt{git push -u origin main}.

% ------------------------------------------------------------------ B
\section{Resultados econométricos completos}

\subsection{Cointegración de Johansen (sensibilidad a muestra y determinística)}

\begin{table}[htb]\centering\small
\caption{Prueba de la traza de Johansen, $H_0\!: r=0$ vs.\ $r\le1$
($k=1$ rezago en diferencias).}
\begin{tabular}{llcccc}
\toprule
Determinística & Muestra & Traza $r{=}0$ & VC 90\,\% & VC 95\,\% & ¿Rechaza 95\,\%?\\
\midrule
Constante            & 2000-2025 & 3,87  & 13,43 & 15,49 & No\\
Constante            & 2000-2019 & 4,64  & 13,43 & 15,49 & No\\
Constante+tendencia  & 2000-2025 & 10,57 & 16,16 & 18,40 & No\\
Constante+tendencia  & 2000-2019 & \textbf{19,88} & 16,16 & 18,40 & \textbf{Sí}\\
\bottomrule
\end{tabular}
\end{table}

Complemento de Engle-Granger (2000-2025): regresión estática
$\ln D_t = 1{,}030 + 0{,}742\,\ln Y_t$ ($R^2=0{,}985$); ADF sobre los
residuos $=-2{,}55$ ($p\approx0{,}0105$): se rechaza no cointegración al
5\,\%. La evidencia conjunta respalda una relación de largo plazo, cuya
inferencia se debilita con el quiebre de 2020, el mismo mecanismo que el BtM
reproducirá de forma permanente sobre la demanda medida.

\subsection{Combinación de pronósticos}

Ponderaciones por inverso del error cuadrático medio de pronósticos a un paso
en evaluación pseudo-fuera-de-muestra 2015-2025 (reestimación recursiva):
VEC 0,456; VAR(1) en diferencias 0,544. El procedimiento replica, a escala,
el esquema oficial (Tabla~1 del informe). Proyección resultante del SIN
orgánico y demanda bruta (con cuña incremental GCE+ME de la UPME):

\begin{table}[htb]\centering\small
\caption{Proyección propia 2026-2040 (GWh-año, escenario medio).}
\begin{tabular}{lrrrrr}
\toprule
 & 2026 & 2030 & 2035 & 2038 & 2040\\
\midrule
SIN orgánico (modelo)        & 86.879 & 98.248 & 114.569 & 125.633 & 133.595\\
Cuña incremental GCE+ME      & 1.829  & 3.661  & 7.820   & 11.520  & 13.987\\
\textbf{Demanda bruta}       & \textbf{88.708} & \textbf{101.909} & \textbf{122.389} & \textbf{137.153} & \textbf{147.582}\\
Referencia UPME jul-2024 (SIN+GCE+ME+GD) & 89.253 & 97.682 & 111.295 & 121.569 & n.d.\\
\bottomrule
\end{tabular}
\end{table}

La senda propia crece más rápido que la referencia oficial (TCAC 3,8\,\% vs.\
2,99\,\% de la rev.\ jul-2026) por dos razones: la elasticidad estimada en la
muestra completa (0,75) recoge el ciclo expansivo 2021-2025, y la cuña GCE se
toma del reporte oficial sin depurar probabilidades de materialización. Para
efectos del análisis BtM lo relevante es el \emph{diferencial} bruta-neta,
que es robusto al nivel de la línea base.

\subsection{Marco de validación: resultados completos}

\begin{table}[htb]\centering\small
\caption{V1. Backtest retrospectivo en Colombia: pronóstico condicional al
PIB observado, por corte de entrenamiento y modelo.}
\begin{tabular}{clcc}
\toprule
Entrenamiento & Modelo & MAPE (\%) & RMSE (GWh)\\
\midrule
\multirow{3}{*}{2000-2019 (pron.\ 2020-2025)}
 & Combinado VEC+VAR & 2,15 & 2.043\\
 & Caminata aleatoria c/deriva & 2,02 & 1.835\\
 & AR(1) univariado & 2,47 & 2.097\\
\midrule
\multirow{3}{*}{2000-2021 (pron.\ 2022-2025)}
 & Combinado VEC+VAR & \textbf{1,15} & \textbf{1.235}\\
 & Caminata aleatoria c/deriva & 1,79 & 1.552\\
 & AR(1) univariado & 2,36 & 1.983\\
\bottomrule
\end{tabular}
\end{table}

\begin{table}[htb]\centering\small
\caption{V2. España: validación del estimador de reconstrucción del
Paso~1. Horas equivalentes calibradas únicamente con 2025
($h=1.174$~h-año); generación estimada
$\widehat{G}_t=\bar{P}_t\cdot h$ vs.\ real (APPA).}
\begin{tabular}{lrrr}
\toprule
Año & $\widehat{G}$ estimada (GWh) & $G$ real (GWh) & Error (\%)\\
\midrule
2019 & 554   & 925    & $-40{,}1$\\
2020 & 1.095 & 1.656  & $-33{,}9$\\
2021 & 2.100 & 3.007  & $-30{,}2$\\
2022 & 4.442 & 4.564  & $-2{,}7$\\
2023 & 7.261 & 7.262  & $0{,}0$\\
2024 & 9.120 & 9.243  & $-1{,}3$\\
2025 & 10.550 & 10.550 & (calibración)\\
\bottomrule
\end{tabular}
\end{table}

\begin{table}[htb]\centering\small
\caption{V3. Panel Colombia-España en tasas de crecimiento (LSDV,
ec.~4 del informe), errores estándar robustos HC1.}
\begin{tabular}{lc}
\toprule
Parámetro & Valor\\
\midrule
$\hat{c}$ (coeficiente de $\Delta s$ BtM) & $-3{,}079$\\
Error estándar robusto & 1,960\\
IC 95\,\% & $[-6{,}92;\ 0{,}76]$\\
$t$ de $H_0\!:c=-1$ ($p$-valor) & $-1{,}06$ (0,299)\\
Elasticidad PIB de corto plazo, Colombia & 0,226\\
Elasticidad PIB de corto plazo, España (demanda neta) & $-0{,}236$\\
$R^2$ / observaciones & 0,58 / 30\\
\bottomrule
\end{tabular}
\end{table}

Notas metodológicas de V3: en niveles, $s_t$ y $\ln Y_t$ resultan casi
colineales en la corta muestra española (2020-2025) y el coeficiente $c$ no
queda identificado (el ejercicio en niveles arroja signos inestables); por
ello la especificación se planteó en primeras diferencias, donde la variación
de instalación anual $\Delta s$ se separa del ciclo del PIB. La elasticidad
española de corto plazo negativa sobre demanda \emph{neta} es en sí misma
sintomática del enmascaramiento que el término $\Delta s$ corrige. La serie
de participación colombiana usa el registro de capacidad AGPE/GD (235,66~MW
acumulados a 2023 y 339~MW a 2024, según UPME; 645~MW a 2025, rev.\
ene-2026), con generación a factor de planta de 16,5\,\%.

% ------------------------------------------------------------------ C
\section{Escenarios BtM año a año}

\begin{table}[htb]\centering\small
\caption{Capacidad BtM instalada por escenario (MW).}
\begin{tabular}{lrrrr}
\toprule
Año & E0 UPME & E1 CREG 4\,\% & E2 España & E3 PNUMA\\
\midrule
2026 & 896   & 1.669 & 842   & 824\\
2028 & 1.727 & 1.787 & 1.365 & 1.342\\
2030 & 3.331 & 1.918 & 2.119 & 2.188\\
2032 & 3.884 & 2.061 & 3.107 & 3.567\\
2034 & 4.529 & 2.217 & 4.276 & 5.815\\
2035 & 4.891 & 2.303 & 4.903 & 7.424\\
2036 & 5.282 & 2.392 & 5.539 & 7.424\\
2038 & 6.159 & 2.581 & 6.794 & 7.424\\
2040 & 7.182 & 2.777 & 7.964 & 7.424\\
\bottomrule
\end{tabular}
\end{table}

\begin{table}[htb]\centering\small
\caption{Demanda neta por escenario (GWh-año) y energía enmascarada en 2040.}
\begin{tabular}{lrrrrr}
\toprule
Trayectoria & 2030 & 2035 & 2040 & TCAC 25-40 & Enmascarado 2040\\
\midrule
Demanda bruta (sin BtM) & 101.909 & 122.389 & 147.582 & 3,81\,\% & 0\\
Neta E0 (tendencial UPME) & 97.094 & 115.320 & 137.201 & 3,31\,\% & 10.381\\
Neta E1 (envolvente CREG) & 99.137 & 119.060 & 143.568 & 3,62\,\% & 4.014\\
Neta E2 (trayectoria España) & 98.846 & 115.302 & 136.071 & 3,25\,\% & 11.511\\
Neta E3 (potencial PNUMA) & 98.746 & 111.658 & 136.851 & 3,29\,\% & 10.731\\
\bottomrule
\end{tabular}
\end{table}

Supuestos: factor de planta FV 16,5\,\% (1~MW $\to$ 1,445~GWh-año;
sensibilidad 15-18\,\% desplaza los cruces del umbral $\pm1$ año);
participación del mercado regulado 68\,\% de la demanda comercial; en E1 se
computa toda la generación BtM contra el umbral (interpretación del cálculo
CREG de 1,2~GW $\leftrightarrow$ 4\,\%); la contabilización alternativa con
excedentes del 40\,\% se reporta en la Tabla~3 del informe.

\begin{figure}[htb]\centering
\includegraphics[width=.8\linewidth]{f3_capacidad_btm.pdf}
\caption{Capacidad BtM por escenario y referencia de España (9,59~GW a 2025).
Elaboración propia.}
\end{figure}

% ------------------------------------------------------------------ D
\section{Figuras complementarias}

\begin{figure}[htb]\centering
\includegraphics[width=.78\linewidth]{f1_demanda_pib.pdf}
\caption{Demanda del SIN y PIB real de Colombia, índice 2000=100. La demanda
crece por debajo del PIB (elasticidad $<1$) y el choque de 2020 rompe
transitoriamente la relación. Fuentes: XM, DANE/Banco Mundial.}
\end{figure}

\begin{figure}[htb]\centering
\includegraphics[width=.78\linewidth]{f6_cointegracion.pdf}
\caption{Relación de cointegración estática $\ln D$-$\ln Y$ (2000-2025).}
\end{figure}

\begin{figure}[htb]\centering
\includegraphics[width=.8\linewidth]{f5_curva_carga.pdf}
\caption{Curva de carga diaria ilustrativa a 2035 bajo E2 (forma típica de
día hábil del SIN, punta nocturna normalizada a la potencia máxima proyectada).
El BtM fotovoltaico deprime el mediodía (${\sim}3{,}7$~GW) sin tocar la punta
de las 19-20~h: el fenómeno de la ``curva de pato''. Elaboración propia.}
\end{figure}

\begin{figure}[htb]\centering
\includegraphics[width=.8\linewidth]{f4_espana_btm.pdf}
\caption{España: potencia BtM acumulada por segmento, 2018-2025 (APPA).
De 302~MW a 9.590~MW en ocho años; el segmento comercial/industrial domina.}
\end{figure}

% ------------------------------------------------------------------ E
\section{Material complementario: regulación y caso España}

\subsection{Marco regulatorio colombiano del BtM}

\begin{itemize}
\item \textbf{Ley 1715 de 2014} (mod.\ Ley 2099 de 2021): promueve FNCER,
autogeneración y entrega de excedentes.
\item \textbf{Res.\ CREG 030 de 2018 $\to$ 174 de 2021}: reglas de AGPE
($\le1$~MW... hasta 5~MW para autogeneración a pequeña escala según ajustes
posteriores) y GD ($<1$~MW en SDL); créditos de energía: excedentes tipo~1
(hasta el consumo) valorados a costo de compra del comercializador, tipo~2 a
precio de bolsa.
\item \textbf{Ley 2294 de 2023 (art.\ 235) y Decreto 2236 de 2023}: crean las
\emph{comunidades energéticas} (AGRC y GDC).
\item \textbf{Doc.\ CREG 901~099 de 2024}: soporte técnico; simula que con
1,2~GW de autogeneración con créditos (equivalente al 4\,\% de la energía del
mercado regulado) el costo unitario de nivel de tensión~1 sube
${\sim}3{,}5$\,\% si todos los excedentes son tipo~1.
\item \textbf{Res.\ CREG 101~072 de 2025 (art.\ 12)}: versión definitiva;
umbral de revisión del 4\,\% de la demanda comercial regulada por mercado de
comercialización; remuneración de excedentes referenciada a MC (precio medio
de contratos del mercado regulado) para estabilizar la señal de precio.
\item \textbf{Contexto de expansión}: el Plan 6GW+ y el balance gremial 2026
(${\sim}4.000$~MW renovables proyectados, de los cuales ${\sim}1.300$~MW en
recursos distribuidos) aceleran la base instalada BtM.
\end{itemize}

\subsection{Indicadores del caso español (APPA 2025 y REE)}

\begin{table}[htb]\centering\small
\caption{España: indicadores BtM 2025.}
\begin{tabular}{lr}
\toprule
Potencia BtM acumulada & 9.590~MW\\
Nueva potencia 2025 (residencial / C\&I) & 1.214~MW (368 / 846)\\
Generación BtM aprovechada 2025 & 10.550~GWh\\
Generación no aprovechada (vertidos limitados) & 2.183~GWh\\
Cobertura de la demanda nacional (256.085~GWh) & 4,1\,\%\\
Punta anual del sistema (15-ene-2025, 20:57) & 40.070~MW\\
Almacenamiento BtM instalado en 2025 & 339~MWh ($+119$\,\%)\\
``Demanda + autoconsumo'' publicada por REE (2025) & 269.753~GWh ($+3{,}7$\,\%)\\
\bottomrule
\end{tabular}
\end{table}

Series anuales empleadas en la validación (APPA, figuras 8, 10 y 11 del
informe): producción de autoconsumo aprovechada 2018-2025 de 446, 925,
1.656, 3.007, 4.564, 7.262, 9.243 y 10.550~GWh; demanda nacional b.c.\
2020-2025 de 250.051, 256.546, 250.471, 246.665, 249.000 y 256.085~GWh; y
energía no aprovechada 2018-2025 de 120, 249, 431, 749, 1.384, 1.899, 2.094
y 2.183~GWh.

Lecciones para Colombia: (i)~la adopción es no lineal y dominada por el
segmento comercial-industrial; (ii)~la cobertura energética del 4\,\% se
alcanzó en ${\sim}7$ años desde el despegue regulatorio (RD 244/2019);
(iii)~la punta del sistema no se redujo, pues el valor de capacidad del BtM
fotovoltaico es bajo sin almacenamiento; (iv)~el evento de ``cero
energético'' de abril de 2025 disparó el almacenamiento detrás del contador
(+119\,\%), añadiendo un servicio de resiliencia que la regulación colombiana
aún no valora; y (v)~la falta de un registro oficial completo de
instalaciones es la principal debilidad estadística señalada por el sector;
Colombia puede evitarla fortaleciendo desde ya los reportes de la
Res.~CREG 174 de 2021 y el registro AGPE de la UPME, insumo directo del
Paso~1 de la metodología propuesta.

\subsection{Sobre los datos}

Demanda SIN 2017-2025 verificada contra XM (2024: 82.084,9~GWh; 2025
estimado con el crecimiento oficial de 2,62\,\%); 2000-2016 compilada de los
informes anuales de XM y regenerable con \texttt{descargar\_datos.py}
(API pydataxm). PIB real de Colombia y España en moneda constante de 2015
del Banco Mundial (serie NY.GDP.MKTP.KN, consistente con DANE e INE). La
proyección oficial UPME jul-2024 se transcribió del documento del curso; de
la rev.\ jul-2026 se tomaron las tasas de crecimiento y las anclas de GD
(3.331~MW a 2030 y 7.182~MW a 2040) publicadas por la UPME. Las series
españolas de la validación provienen del informe APPA 2025 del material del
curso y de REE (nota de prensa del sistema eléctrico 2025, que introduce la
serie oficial ``demanda + autoconsumo'': 269.753~GWh).

% ------------------------------------------------------------------ F
\section{Código fuente principal (extracto)}

Extracto de \texttt{src/ajuste\_btm.py} con la construcción de los
escenarios; el repositorio contiene el código completo comentado.

\begingroup\footnotesize
\begin{verbatim}
GWH_POR_MW = 8760 * FACTOR_PLANTA_FV / 1000.0   # ~1,45 GWh-anio por MW

def esc_tendencial_upme():
    """E0: anclas UPME rev. jul-2026 (645 MW en 2025, rev. ene-2026)."""
    return interp_geometrica({2025: 645, 2030: 3331, 2040: 7182}).loc[H]

def esc_creg_4pct(bruta):
    """E1: capacidad cuya generacion anual equivale al 4 % de la
    demanda comercial regulada (interpretacion conservadora art. 12)."""
    mr = PART_REGULADO * bruta["demanda_bruta_gwh"]
    return (UMBRAL_CREG * mr / GWH_POR_MW)

def esc_espana(bruta):
    """E2: la participacion BtM sobre demanda replica la senda
    espanola: 4,1 % al septimo anio del despegue; logistica s_max=9 %."""
    s_max, s_2025, s_2032 = 0.09, 0.011, 0.041
    def inv(s): return np.log(s_max / s - 1.0)
    k  = (inv(s_2025) - inv(s_2032)) / (2032 - 2025)
    t0 = 2025 + inv(s_2025) / k
    s  = s_max / (1.0 + np.exp(-k * (np.array(H) - t0)))
    return pd.Series(s * bruta["demanda_bruta_gwh"].values / GWH_POR_MW,
                     index=H)

def esc_pnuma():
    """E3: potencial de mercado PNUMA (7.424 MWp) alcanzado en 2035."""
    return interp_geometrica({2025: 645, 2035: 7424, 2040: 7424}).loc[H]

# demanda neta y umbral CREG
for e in ["E0_upme", "E1_creg4", "E2_espana", "E3_pnuma"]:
    gen = esc[f"{e}_mw"] * GWH_POR_MW
    res[f"{e}_neta_gwh"] = bruta["demanda_bruta_gwh"] - gen
    res[f"{e}_pct_umbral_total"] = 100 * gen / mr             # 100 % entregado
    res[f"{e}_pct_umbral_fe"] = 100 * 0.40 * gen / mr         # excedentes 40 %
\end{verbatim}
\endgroup
